\section{Lecture 3: Compactness, Free Condensed Abelian Groups, and Cohomology}
\label{lec:3}

Scholze's second lecture continues the foundations. It opens by recalling the
site of light profinite sets and the crucial feature of its Grothendieck
topology --- that sequential limits of covers are again covers --- and by
revisiting the topological realization functor, improving the statement of
\cref{lec:2} to the identification of \emph{metrizably compactly generated} with
\emph{sequential}. It then isolates the topos-theoretic notions of
\emph{quasicompact} and \emph{quasiseparated} object, checks that they recover
the expected classes of spaces, and uses them to explain why the formalism is
forced to use totally disconnected test objects and a finitary topology with the
Cantor set available. The bulk of the lecture develops light condensed abelian
groups: the tensor product and the free-abelian-group functor (with the running
example $\mathbb{Z}[\mathbb{R}]$), and the three good properties of the light
setting --- exactness of countable products, exactness of sequential limits of
surjections, and internal projectivity of the free group on a convergent
sequence --- the last proved in detail. The lecture closes with a comparison to
all condensed abelian groups, the AB axioms, and the definition of cohomology in
the condensed topos.

\subsection{Recollections: the site of light profinite sets}

Recall from \cref{lec:2} the category of light profinite sets, described
equivalently as
\[
\begin{aligned}
\ProFin_{\mathbb{N}}(\mathrm{Fin})
&\ \simeq\ \{\text{metrizable totally disconnected compact Hausdorff spaces}\} \\
&\ =\ \{\text{countable Boolean algebras}\}^{\mathrm{op}},
\end{aligned}
\]
i.e.\ as sequential (countable) pro-systems of finite sets, as metrizable
totally disconnected compact Hausdorff spaces, or dually as countable Boolean
algebras. We equip this category with the Grothendieck topology whose covers are
generated by two families:
\begin{itemize}
\item finite disjoint unions, and
\item \emph{all} surjective maps.
\end{itemize}
As stressed in \cref{lec:2}, this has the important consequence, which Scholze
underlined again, that
\begin{equation}
\label{eq:3:sequential-limits-of-covers}
\text{sequential limits of covers are still covers}
\end{equation}
(cf.\ \cref{prop:2:countable-limits}: a sequential limit of surjections of light
profinite sets is surjective). This property is what will make the resulting
homological algebra of light condensed abelian groups behave well, and Scholze
returns to it repeatedly below.

Light condensed sets are the sheaves for this topology,
\[
\CondSetl=\mathrm{Shv}\bigl(\ProFin_{\mathbb{N}}(\mathrm{Fin})\bigr).
\]
Whenever one has sheaves on a subcanonical site there is a Yoneda embedding of
the generating category, so the light profinite sets sit inside $\CondSetl$ via
\begin{equation*}
S\ \longmapsto\ \bigl(T\mapsto \Homop_{\ProFin_{\mathbb{N}}(\mathrm{Fin})}(T,S)
=\Cont(T,S)\bigr),
\end{equation*}
sending a profinite set $S$ to the sheaf of (continuous) maps into it. A general
feature is that the image of the Yoneda embedding generates $\CondSetl$ under
colimits: every light condensed set is built from light profinite sets by a
gluing procedure.

There is also the comparison to topological spaces of \cref{lec:2}: a
topological space $A$ gives the light condensed set
\begin{equation*}
\underline{A}\colon S\longmapsto \Cont(S,A),
\end{equation*}
and this functor $A\mapsto\underline{A}$ has a left adjoint
$X\mapsto X(*)_{\mathrm{top}}$, the underlying set $X(*)$ equipped with the
quotient topology from the surjection $\coprod_{\alpha\in X(\text{Cantor})}
(\text{Cantor set})\to X(*)$ (\cref{prop:2:top-realization}).

\subsection{Topological realization: sequential spaces}

Scholze improved on one point from \cref{lec:2}. In forming the quotient
topology on $X(*)_{\mathrm{top}}$ one may use, in place of Cantor sets, the
convergent sequences: crediting a remark of Morita after the previous lecture,
the coproduct
\begin{equation*}
\coprod_{\beta\in X(\mathbb{N}\cup\{\infty\})}(\mathbb{N}\cup\{\infty\})
\ \longrightarrow\ X(*)_{\mathrm{top}}
\end{equation*}
is also a quotient map of topological spaces. The reason is that the convergence
of the Cantor set can be detected by sequences; more precisely, mapping all
convergent sequences to the Cantor set is still a quotient map, so that in
$\mathrm{Top}$ one may write
\begin{equation}
\label{eq:3:cantor-colim}
\text{Cantor set}\ \simeq\ \operatorname*{colim}\ (\text{countable closed subsets})
\end{equation}
as a colimit over its countable closed subsets, e.g.\ finite unions of convergent
sequences. Scholze emphasized that \eqref{eq:3:cantor-colim} holds \emph{in
$\mathrm{Top}$ but not in condensed sets} --- a point he returns to below.

Consequently the corresponding topological space is determined by its convergent
sequences alone, and the notion of being \emph{metrizably compactly generated}
from \cref{lec:2} is the same as being \emph{sequential}: continuity may be
checked on convergent sequences (a map is continuous iff it sends convergent
sequences to convergent sequences).

\begin{proposition}[Sequential spaces embed]
\label{prop:3:sequential-embed}
The functor $A\mapsto\underline{A}$ restricts to a fully faithful embedding
\begin{equation*}
\{\text{sequential topological spaces}\}\ \hookrightarrow\ \CondSetl.
\end{equation*}
\end{proposition}

\begin{remark}[Sequential versus first-countable]
\label{rmk:3:sequential-first-countable}
Scholze was unsure off-hand how ``sequential'' compares with ``first
countable''; in fact every first-countable space is sequential and the
containment is strict. (The large discrete space floated in the discussion does
not separate the two, since discrete spaces are always first countable.) He
noted that in an earlier version of the theory the relevant condition was
first-countability with all points closed; here one no longer needs all points
closed, having dropped the set-theoretic bounds.
\end{remark}

\subsection{Quasicompact and quasiseparated objects}

Scholze now recalled why the Cantor set is allowed as a test object, through two
general topos-theoretic notions. In any topos --- sheaves on any site --- there
is an intrinsic notion of an object being ``compact'' and of being
``Hausdorff''.

\begin{definition}[Quasicompact, quasiseparated]
\label{def:3:qc-qs}
Let $X$ be a light condensed set.
\begin{enumerate}
\item $X$ is \emph{quasicompact} (qc) if any cover
$\coprod_i X_i\twoheadrightarrow X$ admits a finite subcover. For light condensed
sets this amounts to: there exists a surjection from the Cantor set,
\begin{equation*}
\{0,1\}^{\mathbb{N}}\twoheadrightarrow X,
\end{equation*}
or $X=\emptyset$.
\item $X$ is \emph{quasiseparated} (qs) if for all quasicompact $Y,Z\to X$ the
fiber product $Y\times_X Z$ is quasicompact. For light condensed sets this
amounts to: for all pairs of maps $\{0,1\}^{\mathbb{N}}\to X$ from the Cantor
set,
\begin{equation}
\label{eq:3:qs}
\{0,1\}^{\mathbb{N}}\times_X\{0,1\}^{\mathbb{N}}\quad\text{is quasicompact.}
\end{equation}
\end{enumerate}
If both hold, $X$ is called \emph{qcqs}.
\end{definition}

\begin{remark}[Terminology and provenance]
\label{rmk:3:grothendieck-qs}
The condition of being quasiseparated is Hausdorffness at the level of the
topos. It was introduced by Grothendieck in SGA~4; using the algebraic-geometry
convention where ``separated'' means Hausdorff, and this being a general
topos-theoretic notion rather than a specific separatedness in algebraic
geometry, he called it \emph{quasiseparated}. The intuition for
\eqref{eq:3:qs}: a map $\{0,1\}^{\mathbb{N}}\to X$ need not be injective, so it
factors over a quotient of the Cantor set; quasiseparatedness declares that this
is the quotient by a \emph{closed} equivalence relation on the Cantor set, one
way to express Hausdorffness.
\end{remark}

For a finitary topology the quasicompact objects are exactly those covered by a
finite union of generating objects, such as $\{0,1\}^{\mathbb{N}}$. Scholze used
this to explain the design of the site:

\begin{remark}[Why the Cantor set must be a test object]
\label{rmk:3:need-cantor}
If one allowed only the convergent sequence $\mathbb{N}\cup\{\infty\}$ in the
test category, then the quasicompact objects would all be countable, and in
particular the Cantor set would \emph{not} be quasicompact --- one could only
access it as a huge colimit of smaller objects, as in \eqref{eq:3:cantor-colim}.
Since the Cantor set is morally a basic compact object, this would destroy the
intuition attached to quasicompactness. So the abstract topos-theoretic notion
mirrors closely what a compact space should be, and this forces the topology to
be finitary and forces the Cantor set into the formalism.
\end{remark}

The abstract notions recover exactly the expected classes of spaces.

\begin{proposition}[qcqs objects are metrizable compact Hausdorff spaces]
\label{prop:3:qcqs}
Under $A\mapsto\underline{A}$ there is an equivalence
\begin{equation*}
\{\text{qcqs light condensed sets}\}\ \xrightarrow{\ \sim\ }\
\{\text{metrizable compact Hausdorff spaces}\}.
\end{equation*}
If one removes ``light'' on the left, one removes ``metrizable'' on the right.
\end{proposition}

\begin{example}[The interval and its Cantor covering]
\label{ex:3:interval-cantor}
The unit interval $[0,1]$ is a metrizable compact Hausdorff space, hence
quasicompact as a condensed set, so there must be a surjection from the Cantor
set onto it --- and indeed a canonical one. A sequence $(a_0,a_1,a_2,\dots)$ of
$0$'s and $1$'s is sent to the real number with binary expansion
$0.a_0a_1a_2\cdots$, giving
\begin{equation*}
\{0,1\}^{\mathbb{N}}\ \twoheadrightarrow\ [0,1],
\end{equation*}
which is surjective since every point of the interval admits such a
representation. One recovers $[0,1]$ as the quotient by the closed equivalence
relation identifying the two binary expansions of a dyadic rational, e.g.\
$0.0111\cdots=0.1000\cdots$. One genuinely needs something as large as the
Cantor set to surject onto $[0,1]$.
\end{example}

Dropping the compactness but keeping the Hausdorff condition, the quasiseparated
objects also have a standard description.

\begin{proposition}[Quasiseparated light condensed sets]
\label{prop:3:qs-ind}
There is an equivalence
\begin{equation*}
\{\text{qs light condensed sets}\}\ \simeq\
\Ind_{\mathrm{inj}}\bigl(\{\text{metrizable compact Hausdorff spaces}\}\bigr),
\end{equation*}
the ind-category of metrizable compact Hausdorff spaces along injective
transition maps. Concretely the objects are filtered systems
$\{X_i\}_{i\in I}$ ($I$ a filtered poset) of metrizable compact Hausdorff spaces
in which all transition maps are injective --- hence automatically closed
immersions --- and one should think of such an object as a rising union of its
quasicompact subspaces $X_i$.
\end{proposition}

\begin{remark}[Comparison with compactly generated weak Hausdorff spaces]
\label{rmk:3:cgwh}
Inside the quasiseparated light condensed sets one finds the
\emph{compactly generated weak Hausdorff spaces} (``CGWH'') of algebraic
topology, which by the discussion above are exactly the sequential ones. The two
categories are nevertheless genuinely different. In $\mathrm{Top}$ (or CGWH) the
Cantor set is the colimit \eqref{eq:3:cantor-colim} of its countable closed
subsets, e.g.\ finite unions of convergent sequences; but this colimit
identity \emph{fails} in quasiseparated light condensed sets. There the Cantor
set is a single quasicompact object, whereas the colimit of countable closed
subsets is a genuine (formal) filtered colimit --- so the two are very different
as condensed sets, the latter not even being quasicompact. Scholze noted that
the discrepancy only appears once the indexing colimit is large: restricting to
countable systems, the two notions agree. Removing ``light'' throughout removes
``metrizable''.
\end{remark}

\subsection{Light condensed abelian groups}

We now put algebra on top of the formalism. Light condensed abelian groups admit
the two equivalent descriptions
\begin{equation*}
\CondAbl=\{\text{abelian group objects in }\CondSetl\}
=\mathrm{Shv}\bigl(\ProFin_{\mathbb{N}}(\mathrm{Fin}),\,\mathrm{Ab}\bigr),
\end{equation*}
either abelian group objects in light condensed sets, or sheaves of abelian
groups on the site; the same remark applies to any algebraic structure (e.g.\
condensed rings are ring objects in $\CondSetl$, equivalently sheaves of rings).

From the general theory of sheaves of abelian groups on a site (cf.\
\cref{prop:2:grothendieck}), $\CondAbl$ is a Grothendieck abelian category: all
limits and colimits exist and filtered colimits are exact. It also carries a
symmetric monoidal tensor product, with the following features.

\begin{notation}[Unit and tensor product]
\label{not:3:tensor}
The unit object is $\underline{\mathbb{Z}}$, the (discrete) condensed abelian
group associated to $\mathbb{Z}$, with
\begin{equation*}
\underline{\mathbb{Z}}\colon S\longmapsto \Cont(S,\mathbb{Z})
=\{\text{locally constant maps }S\to\mathbb{Z}\}.
\end{equation*}
The tensor product $M\otimes N$ is the sheafification of the presheaf
\begin{equation*}
S\longmapsto M(S)\otimes N(S).
\end{equation*}
\end{notation}

The forgetful functor $\CondAbl\to\CondSetl$ has a left adjoint, the free
condensed abelian group functor.

\begin{definition}[Free condensed abelian group]
\label{def:3:free-group}
The free condensed abelian group on a light condensed set $X$ is
$\mathbb{Z}[X]$, the sheafification of the presheaf
\begin{equation*}
S\longmapsto \mathbb{Z}[X(S)],
\end{equation*}
the levelwise free abelian group. Its underlying group is the ordinary free
abelian group $\mathbb{Z}[X(*)]$ on the underlying set; the construction puts a
condensed (``topological'') structure on it.
\end{definition}

This is a structure one rarely contemplates for topological spaces --- one does
not usually form the free abelian group on a topological space --- but here it
exists and is completely fundamental.

\begin{example}[The free group on the reals]
\label{ex:3:ZR}
Take $X=\mathbb{R}$ (from now on the underline is usually dropped, everything
having implicitly become a condensed set) and form $\mathbb{Z}[\mathbb{R}]$. Its
elements are finite formal sums
\begin{equation*}
\sum_{x\in\mathbb{R}} n_x[x],\qquad n_x\in\mathbb{Z}\text{ almost all }0,
\end{equation*}
thought of as finite integer combinations of Dirac measures, but now carrying a
topology recording that a point $x$ may move continuously: $[x]-[y]$ is a nonzero
element in general, yet it collapses to $0$ as $y\to x$. It is a little awkward
to describe the open subsets of $\mathbb{Z}[\mathbb{R}]$ as a topological group,
but as a condensed set it is transparent. Since the free construction commutes
with colimits,
\begin{equation*}
\mathbb{Z}[\mathbb{R}]
=\operatorname*{colim}_{c}\ \mathbb{Z}\bigl[[-c,c]\bigr],
\end{equation*}
the colimit over intervals $[-c,c]$, now with compact pieces, and each
$\mathbb{Z}[I]$ ($I$ a compact interval) is the rising union
\begin{equation*}
\mathbb{Z}[I]=\bigcup_{n\in\mathbb{N}}\mathbb{Z}[I]_{\leq n},
\qquad
\mathbb{Z}[I]_{\leq n}=\Bigl\{\textstyle\sum_x n_x[x]\ :\ \sum_x|n_x|\leq n\Bigr\},
\end{equation*}
each $\mathbb{Z}[I]_{\leq n}$ being compact Hausdorff. Indeed, whenever $K$ is a
compact Hausdorff space, the integer combinations of points of $K$ with
$\sum|n_x|\leq n$ carry a canonical compact Hausdorff topology; the free
construction produces it automatically. Passing from a profinite set into
$\mathbb{Z}[\mathbb{R}]$, one always lands in one of these compact subsets.
\end{example}

\begin{remark}[Nontrivial identifications require surjective covers]
\label{rmk:3:collapse-needs-surjections}
The subtle content of \cref{ex:3:ZR} is precisely the identification collapsing
$[x]-[y]$ to $0$ as $y\to x$. Making this come out correctly uses that the
Grothendieck topology allows not only finite disjoint unions but also surjective
maps as covers; with the smaller topology the free construction would give the
wrong answer.
\end{remark}

\begin{remark}[Free groups on group objects are rings; real powers]
\label{rmk:3:ZG-ring}
If $X$ already carries a group structure, then $\mathbb{Z}[X]$ acquires a ring
structure, the multiplication coming from the group law on $X$; so
$\mathbb{Z}[\mathbb{R}]$ is naturally a condensed ring. Regarding $\mathbb{R}$
multiplicatively as the group of exponents of a formal variable $T$, this ring
is $\mathbb{Z}[T^{\mathbb{R}}]$: it adjoins to $\mathbb{Z}$ a variable together
with all of its real powers. This is a first incarnation of the object
$K\langle T^{\mathbb{R}}\rangle$ that motivated \cref{lec:2}
(see~\eqref{eq:2:global-idea}), where its very nature was deliberately left open
--- only the wish that the condensed formalism should make such objects at least
\emph{exist} as condensed rings; the free construction now produces one for
free. Editorially: this is not yet the completed or topologized version one
ultimately wants.
\end{remark}

\subsection{Three properties special to the light setting}

We turn to features specific to \emph{light} condensed abelian groups.

\begin{theorem}[Good properties of $\CondAbl$]
\label{thm:3:three-properties}
In $\CondAbl$:
\begin{enumerate}
\item countable products are exact;
\item sequential limits of surjective maps are surjective; and
\item the free group $\mathbb{Z}[\mathbb{N}\cup\{\infty\}]$ on the convergent
sequence is \emph{internally projective}.
\end{enumerate}
Properties \textup{(1)} and \textup{(2)} in fact hold in all condensed abelian
groups (there \emph{all} products are exact); property \textup{(3)} is special to
the light setting.
\end{theorem}

Property (2) is what one repeatedly needs in functional analysis: sequential
limits of surjections --- e.g.\ a Fr\'echet space as a sequential limit of Banach
spaces --- should behave well.

\begin{definition}[Internal projectivity]
\label{def:3:internally-projective}
Recall that in any Grothendieck abelian category $P$ is \emph{projective} if
$\Ext^i(P,-)=0$ for $i>0$. When the category is symmetric monoidal one has an
\emph{internal} $\Ext$: for $M,N\in\CondAbl$,
\begin{equation*}
\underline{\Ext}^i(M,N):=\text{sheafification of}\quad
S\longmapsto \Ext^i\bigl(M\otimes\mathbb{Z}[S],\,N\bigr).
\end{equation*}
An object $P$ is \emph{internally projective} if $\underline{\Ext}^i(P,-)=0$ for
$i>0$.
\end{definition}

\begin{proof}[Proof of \textup{(1)} and \textup{(2)}]
\emph{Reduction of \textup{(1)} to \textup{(2)}.} For a countable product of
maps, products of injections are injections and preserve exactness on the left;
the only point needing proof is that a countable product of surjections
$f_n\colon M_n\twoheadrightarrow N_n$ is surjective. Write the map
$\prod_n M_n\to\prod_n N_n$ as the sequential limit
\begin{equation*}
\prod_n M_n\to\prod_n N_n\ =\
\varprojlim_m\ \Bigl(\ \prod_{n\leq m} M_n\times\prod_{n>m} N_n\
\twoheadrightarrow\ \prod_n N_n\ \Bigr).
\end{equation*}
Each map in the limit is surjective: finite products are finite direct sums,
which preserve surjections and are exact, so one can lift the first $m$
coordinates while leaving the rest. Thus $\prod f_n$ is a sequential limit of
surjections, and \textup{(2)} finishes the argument.

\emph{Proof of \textup{(2)}.} Let $M_\infty=\varprojlim(\cdots\to M_2\to
M_1\twoheadrightarrow M_0)$ be a sequential limit of surjections of light
condensed abelian groups; we show $M_\infty\to M_0$ is surjective in the sheaf
sense. Surjectivity means: for every light profinite set $S$ with a map $S\to
M_0$, there is a surjection $S_\infty\twoheadrightarrow S$ from a light profinite
set and a lift $S_\infty\to M_\infty$. Since $M_1\twoheadrightarrow M_0$ is
surjective, there is a light profinite set $S_1\twoheadrightarrow S$ with a lift
$S_1\to M_1$; inductively there is $S_{i+1}\twoheadrightarrow S_i$ with a lift to
$M_{i+1}$. Put
\begin{equation*}
S_\infty=\varprojlim_i S_i\ \in\ \ProFin_{\mathbb{N}}(\mathrm{Fin}).
\end{equation*}
This is a countable limit of surjections, hence still surjective onto $S$
(\cref{prop:2:countable-limits}), so $S_\infty\twoheadrightarrow S$ is a cover;
and the compatible lifts assemble to a map $S_\infty\to M_\infty$. This is
exactly the required lift, and \eqref{eq:3:sequential-limits-of-covers} --- that
sequential limits of covers are covers --- is the crux.
\end{proof}

\begin{remark}[Why totally disconnected building blocks]
\label{rmk:3:totally-disconnected-forced}
Property (2), i.e.\ that countable limits of covers are covers, is exactly what
forces the test objects to be totally disconnected. Suppose one tried instead a
``smooth'' site of manifolds with the usual topology of open covers. Cover a
small interval by two subintervals, cover each of those by two smaller ones, and
continue; in the limit the intervals shrink to a cloud of dust --- a Cantor set
--- with no reasonable smooth structure. So iterating covers drives one to
totally disconnected compact Hausdorff spaces as building blocks, and among these
Scholze takes the metrizable ones with all surjections as covers.
\end{remark}

\begin{remark}[Provenance: the pro-\'etale topology]
\label{rmk:3:proetale-origin}
Scholze noted, as a historical aside, that this prototype is not new: the
condensed formalism grows out of the pro-\'etale topology for schemes
\cite{bhatt2014proetaletopologyschemes}, where the very wish was that limits of
surjective maps remain surjective, so that certain sheaves arising as inverse
limits would be well behaved. Allowing countable (or all) limits of covers to
remain covers is the origin of the whole theory. He added that any surjection of
light profinite sets can be written as a sequential limit of open covers --- for
instance $\mathbb{N}\cup\{\infty\}$ is covered by two copies of
$\mathbb{N}\cup\{\infty\}$ meeting at $\infty$ --- so if one wants open covers
and their sequential limits, one must allow all surjective maps as covers.
\end{remark}

\subsection{Internal projectivity of the free group on a convergent sequence}

Before the proof, a warning: the phenomenon appears only after passing to
abelian groups.

\begin{remark}[$\mathbb{N}\cup\{\infty\}$ is not projective as a condensed set]
\label{rmk:3:seq-not-projective}
The convergent sequence $\mathbb{N}\cup\{\infty\}$ is \emph{not} projective in
$\CondSetl$. Take the surjection of light profinite sets
\begin{equation*}
S'=(2\mathbb{N}\cup\{\infty\})\sqcup\bigl((2\mathbb{N}+1)\cup\{\infty\}\bigr)
\ \twoheadrightarrow\ S=\mathbb{N}\cup\{\infty\},
\end{equation*}
a disjoint union of two convergent sequences (the evens, the odds), each with its
own limit point, both mapping to $\infty$. The identity map
$\mathbb{N}\cup\{\infty\}\to S$ admits no lift to $S'$: an integer has a unique
preimage, so a lift would send the even and odd subsequences to the two
\emph{different} limit points of $S'$, and the resulting map does not converge.
% board 16 @ 01:11:30 — /home/deven/documents/coding/vms/gpu/home/notetaker/output/peter-scholze-324-analytic-stacks/boards/board-16.jpg
\[
\begin{tikzcd}
& S'=(2\mathbb{N}\cup\{\infty\})\sqcup\bigl((2\mathbb{N}+1)\cup\{\infty\}\bigr) \arrow[d, two heads] \\
\mathbb{N}\cup\{\infty\} \arrow[r] \arrow[ur, dashed, "\not\exists"] & S=\mathbb{N}\cup\{\infty\}
\end{tikzcd}
\]
Miraculously, once one passes to abelian groups this obstruction disappears: any
null sequence \emph{can} be lifted, which is the content of
\cref{thm:3:three-properties}(3).
\end{remark}

It is convenient to prove the statement for a direct factor of
$\mathbb{Z}[\mathbb{N}\cup\{\infty\}]$.

\begin{definition}[The free group on a null sequence]
\label{def:3:null-sequence-group}
Set
\begin{equation*}
M:=\mathbb{Z}[\mathbb{N}\cup\{\infty\}]\big/\mathbb{Z}\cdot[\infty],
\end{equation*}
the free condensed abelian group on the pointed convergent sequence (the limit
point sent to $0$). It is the free group on a \emph{null sequence} and classifies
null sequences: maps out of $M$ are null sequences in the target. Since
$\{\infty\}\hookrightarrow\mathbb{N}\cup\{\infty\}$ splits (project to a point),
$M$ is a direct factor of $\mathbb{Z}[\mathbb{N}\cup\{\infty\}]$, the complementary
factor $\mathbb{Z}$ being projective; so it suffices to prove $M$ internally
projective.
\end{definition}

\begin{proof}[Proof of \cref{thm:3:three-properties}(3)]
We prove ordinary projectivity of $M$ and remark afterward on the internal
statement. Let $\widetilde{N}\twoheadrightarrow N$ be a surjection of light
condensed abelian groups and $M\to N$ a map --- equivalently a null sequence in
$N$, i.e.\ a map $\mathbb{N}\cup\{\infty\}\to N$ sending $\infty\mapsto 0$. We
must lift it to $\widetilde{N}$.
% board 16 @ 01:11:30 — /home/deven/documents/coding/vms/gpu/home/notetaker/output/peter-scholze-324-analytic-stacks/boards/board-16.jpg
\[
\begin{tikzcd}
& \widetilde{N} \arrow[d, two heads] \\
M \arrow[r] \arrow[ur, dashed, "\exists?"] & N
\end{tikzcd}
\]
Since $\widetilde{N}\twoheadrightarrow N$ is a surjection of condensed abelian
groups, the null sequence lifts after a cover: there is a surjection
$S\twoheadrightarrow\mathbb{N}\cup\{\infty\}$ from a light profinite set together
with a map $g\colon S\to\widetilde{N}$ lifting it. Shrinking $S$, we may assume
that over each integer the fibre is a single point, i.e.
\begin{equation*}
S\times_{\mathbb{N}\cup\{\infty\}}\mathbb{N}\ \xrightarrow{\ \sim\ }\ \mathbb{N},
\end{equation*}
so that $S$ is merely some compactification of $\mathbb{N}$. Let
\begin{equation*}
S_\infty:=S\times_{\mathbb{N}\cup\{\infty\}}\{\infty\}\ \subseteq\ S,
\end{equation*}
the fibre over $\infty$, a closed subspace.

Here lightness enters decisively: $S_\infty$ is light, hence injective in
$\ProFin(\mathrm{Fin})$ (\cref{prop:2:injective}), so the inclusion
$S_\infty\hookrightarrow S$ admits a retraction $r\colon S\to S_\infty$. Consider
the map
\begin{equation*}
g-g\circ r\colon\ S\longrightarrow \widetilde{N}.
\end{equation*}
On $S_\infty$ one has $r|_{S_\infty}=\mathrm{id}$, so $g-g\circ r$ vanishes on
$S_\infty$.

Now the square
% board 17 @ 01:13:09 — /home/deven/documents/coding/vms/gpu/home/notetaker/output/peter-scholze-324-analytic-stacks/boards/board-17.jpg
\begin{equation*}
\begin{tikzcd}
S_\infty \arrow[r, hook] \arrow[d] & S \arrow[d, two heads] \\
\{\infty\} \arrow[r, hook] & \mathbb{N}\cup\{\infty\}
\end{tikzcd}
\end{equation*}
is a pushout in light condensed sets --- a form of the gluing condition, valid
because $S\twoheadrightarrow\mathbb{N}\cup\{\infty\}$ is surjective. Descending a
map out of $S$ to $\mathbb{N}\cup\{\infty\}$ requires only that it be constant on
fibres; over the integers the fibres are points, so the sole condition is at the
fibre $S_\infty$ over $\infty$, where indeed $g-g\circ r=0$. Hence $g-g\circ r$
factors through a map
\begin{equation*}
\widetilde{f}\colon \mathbb{N}\cup\{\infty\}\longrightarrow\widetilde{N},
\qquad \widetilde{f}(\infty)=0,
\end{equation*}
i.e.\ a null sequence in $\widetilde{N}$, equivalently a map $M\to\widetilde{N}$.

It remains to check that $\widetilde{f}$ lifts the original map, i.e.\ projects to
it in $N$. Since $\mathbb{Z}[S]\twoheadrightarrow M$ is surjective, it suffices to
check after precomposing with $S$. Projecting $g-g\circ r$ to $N$: the term $g$
projects to the original lift, while the correction $g\circ r$ projects to the
composite $S\to N$ that factors through $r$ and hence through $S_\infty$; because
the original null sequence sends $\infty\mapsto 0$, this composite vanishes on
$S_\infty$ and so projects to $0$ in $N$. Therefore $\widetilde{f}$ projects to
the original null sequence, and $M$ is projective.

For internal projectivity the argument is essentially the same, carrying along an
auxiliary profinite set that one also covers; again the decisive point is that
sequential limits of surjections are surjections.
\end{proof}

\subsection{Comparison with all condensed abelian groups}

It is worth contrasting the light setting with the category $\mathrm{CondAb}$ of
\emph{all} condensed abelian groups.

\begin{remark}[Projectives in $\mathrm{CondAb}$]
\label{rmk:3:condab-projectives}
In $\mathrm{CondAb}$ all products are exact, and there are projective generators
$\mathbb{Z}[S]$ with $S$ extremally disconnected, e.g.\ $S=\beta I$ the
Stone--\v{C}ech compactification of a discrete set $I$. These are huge and
inexplicit --- their structure depends on the model of set theory --- but they
exist by the axiom of choice. The catch is that \emph{none} of them are
internally projective: this involves the badly-behaved products
\begin{equation*}
\beta I\times\beta J\qquad(I,J\ \text{infinite})
\end{equation*}
of extremally disconnected sets. The product $\beta I\times\beta J$ is itself
\emph{not} extremally disconnected, and Scholze suggested this may be what
obstructs internal projectivity. It is a rather severe technical nuisance ---
many arguments really want the projectives to be internally projective.
\end{remark}

\begin{remark}[Projectives and injectives in $\CondAbl$]
\label{rmk:3:light-proj-inj}
By contrast $\CondAbl$ has \emph{no} projective generators: the free group
$\mathbb{Z}[\{0,1\}^{\mathbb{N}}]$ on the Cantor set cannot be covered by a
projective object, so the category does not have enough projectives --- one of
the reasons Scholze was initially hesitant to make the switch, though in practice
it turns out not to matter much. What one does have is the single very explicit
internally projective object $\mathbb{Z}[\mathbb{N}\cup\{\infty\}]$
(equivalently $M$ of \cref{def:3:null-sequence-group}), which is worth a great
deal. Dually, $\CondAbl$ has \emph{enough injectives} --- they exist for general
set-theoretic reasons but cannot be written down explicitly --- whereas
$\mathrm{CondAb}$ has \emph{no} nonzero injective objects at all.
\end{remark}

\begin{remark}[{$\mathbb{Z}[\mathbb{N}\cup\{\infty\}]$ is not projective in
$\mathrm{CondAb}$}]
\label{rmk:3:not-projective-all}
The free group on the convergent sequence, projective in $\CondAbl$, is
\emph{not} projective in $\mathrm{CondAb}$. The universal compactification of
$\mathbb{N}$ is the Stone--\v{C}ech compactification $\beta\mathbb{N}$, and the
surjection $\mathbb{Z}[\beta\mathbb{N}]\to\mathbb{Z}[\mathbb{N}\cup\{\infty\}]$
admits no splitting --- the Stone--\v{C}ech obstruction. So light-projectivity
here is genuinely a light phenomenon, and $\mathbb{Z}[\mathbb{N}\cup\{\infty\}]$
does not even remain projective under the embedding $\CondAbl\hookrightarrow
\mathrm{CondAb}$.
\end{remark}

\begin{remark}[Relation to injective Banach spaces]
\label{rmk:3:banach}
These phenomena are shadows of facts about Banach spaces. Under the duality
whereby a projective object of the condensed world becomes an injective Banach
space, the relevant objects are the spaces $\Cont(S,\mathbb{R})$ with $S$
extremally disconnected (a Stone--\v{C}ech compactification, or a retract of
one); it is a classical theorem that every injective Banach space is a retract of
such a $\Cont(S,\mathbb{R})$. The object dual to the free group on a null
sequence is the Banach space $c_0(\mathbb{N})$ of null sequences, which is
\emph{not} injective but is \emph{separably injective} --- injectivity tested
only against separable Banach spaces \cite{Avil_s_2016}. Scholze remarked that he
arrived at the light-projectivity of $\mathbb{Z}[\mathbb{N}\cup\{\infty\}]$ after
reading the proof that $c_0$ is separably injective. He added that
whether $\Cont(\beta\mathbb{N},\mathbb{R})$ is separably injective was an open
question in the Banach-space literature, and that the analytic/complex-geometry
framework of the course can show it is \emph{not} even separably injective --- so
this space does not behave like a projective object even when tested only against
light condensed abelian groups.
\end{remark}

\begin{remark}[AB axioms]
\label{rmk:3:ab-axioms}
Grothendieck's T\^ohoku axioms for an abelian category are relevant here. The
axiom that all products be exact (AB5*, the dual product form) holds in
$\mathrm{CondAb}$ but fails in $\CondAbl$, where only countable products are
exact. The next axiom, AB6 --- a certain commutation of infinite products with
filtered colimits --- holds whenever there are enough projective generators, in
particular in $\mathrm{CondAb}$, and, restricted to countable products, also in
$\CondAbl$:
\begin{equation*}
\text{AB6 holds for countable products in }\CondAbl.
\end{equation*}
One way to see this: the fully faithful embedding
$\CondAbl\hookrightarrow\mathrm{CondAb}$ commutes with all colimits and with
countable limits, so countable products in $\CondAbl$ agree with those in
$\mathrm{CondAb}$, where AB6 holds; it can also be checked by hand. (This
sharpens \cref{thm:2:condab-properties}(1).)
\end{remark}

\subsection{Cohomology in the condensed topos}

The lecture closed with cohomology. Working in any topos, one automatically
inherits a notion of cohomology, and the point is that in the condensed topos it
recovers the expected invariants of topological spaces even though no geometry
was built into the site.

\begin{definition}[Condensed cohomology]
\label{def:3:cohomology}
Let $X$ be any condensed set and $M$ an abelian group (for the moment a discrete
one; more generally a condensed abelian group). Define
\begin{equation*}
H^{*}(X,M):=\Ext^{*}_{\CondAbl}\bigl(\mathbb{Z}[X],\,M\bigr),
\end{equation*}
the $\Ext$ computed in light condensed abelian groups (the same as in all
condensed abelian groups).
\end{definition}

\begin{theorem}[Recovery of singular cohomology]
\label{thm:3:singular}
If $X$ is a CW complex, then
\begin{equation*}
H^{i}(\underline{X},M)\ \cong\ H^{i}_{\mathrm{sing}}(X,M),
\end{equation*}
the singular cohomology of $X$.
\end{theorem}

\begin{remark}[Geometry from totally disconnected data]
\label{rmk:3:geometry-out-of-nothing}
This is striking: the site used only totally disconnected test objects, yet the
resulting cohomology detects, for instance, that the circle is not contractible.
The interval, as a condensed set, is a genuine interval and not a point; one
passes to homotopy precisely by forming these $\Ext$-groups, since the interval
admits no nonzero map to a discrete abelian group.
\end{remark}

\begin{remark}[{Dold--Thom: $\mathbb{Z}[X]$ models homology}]
\label{rmk:3:dold-thom}
The heuristic behind \cref{thm:3:singular} is the Dold--Thom theorem
\cite{Dold_1958}: the free abelian group $\mathbb{Z}[X]$ --- finite integer
combinations of points of $X$ --- is, once one passes to homotopy, a model for
the homology of $X$. Dualizing (via $\Ext$ into $M$) then gives cohomology,
fitting the picture.
\end{remark}

\begin{remark}[Internal version and continuous group cohomology]
\label{rmk:3:internal-and-group}
Using the internal $\underline{\Ext}$ instead unravels, by adjunction, to
replacing the coefficient $M$ by the continuous functions
$\underline{\Cont}(S,M)$ out of a test object $S$ --- still a condensed abelian
group --- so it does not carry much more information. A more
substantial payoff is for continuous group cohomology: for a topological group,
regarded as a condensed group, it is clear what an action on a condensed abelian
group is and what its cohomology should be, via the general topos-theoretic
answer. This always gives the expected answer, which need not coincide with the
one computed by explicit continuous cochains --- and where it differs, the
condensed answer is the better one. For a CW complex with a local system one gets
a sheaf on condensed sets over $\underline{X}$, and the usual cohomology is
singular sheaf cohomology; the comparison rests on cohomological descent for
surjections of compact Hausdorff spaces.
\end{remark}

\begin{remark}[Why the finitary site is the right one for cohomology]
\label{rmk:3:sequential-topos-bad}
Had one worked instead in the sequential topos generated by
$\mathbb{N}\cup\{\infty\}$ alone, computing the cohomology of the interval would
first require expressing the interval as the huge colimit
\eqref{eq:3:cantor-colim} of countable closed subsets; the cohomology would then
return an enormous derived limit, and it is very unlikely to produce the expected
answer. That the finitary, Cantor-set-based site does give the right answer is a
strong indication that it is the good framework.
\end{remark}

\begin{remark}[Enough projectives, again]
\label{rmk:3:enough-projectives-closing}
Asked whether $\CondAbl$ has enough projective objects, Scholze reiterated that
it does not --- one does not expect a projective cover of the free group on the
Cantor set. In practice one embeds $\CondAbl\hookrightarrow\mathrm{CondAb}$, a
functor preserving all colimits and all countable limits, hence compatible with
many operations, and carries out computations in $\mathrm{CondAb}$, which does
have enough projectives.
\end{remark}
